% Section 8.3, from lectures/L08/appendix/c_power_amplifiers.md.

\section{What a real power amplifier adds}
\label{c8:app:c}

\textbf{Reading, not examinable.} No exercises, nothing in the test suite. Everything in \secref{c8:app:b} still holds; this is the shape of what it leaves out.

\subsection{The power arithmetic, done properly}
\label{c8:sec:c1}
50 W into 8 ohm. Confusing rms with peak here propagates into the rail voltage and everything after it.

\[
  V_{rms} = \sqrt{P R} = \sqrt{50 \times 8} = 20\ \text{V}, \qquad I_{rms} = 2.5\ \text{A}
\]

\[
  V_{peak} = \sqrt{2}\,V_{rms} = 28.3\ \text{V}, \qquad I_{peak} = 3.54\ \text{A}
\]

\textbf{Calling 2.5 A the peak current is the easy mistake,} and it leads to $\pm 20$ V rails, which deliver at most

\[
  P = \frac{V_{rail}^2}{2R} = \frac{400}{16} = 25\ \text{W}
\]

\textbf{half the power the example is about.} A 50 W amplifier into 8 ohm needs $\pm 28.3$ V before any allowance for saturation, the emitter resistors and supply sag: in practice $\pm 32$ V.

\begin{narrowtable}{@{}lll@{}}
  \thead{} & \thead{Idle current} & \thead{Idle dissipation} \\
  \midrule
  Class A & 3.54 A, the full peak & 226 W \\
  Class AB & 120 mA & 7.7 W \\
\end{narrowtable}

\textbf{A factor of 29,} for a stage that measures very nearly as well. That is the whole argument for class AB.

\textbf{Maximum theoretical efficiency,} ideal devices at full output: class A \textbf{25 per cent} resistively loaded and 50 with a current-source load, class B \textbf{78.5 per cent}, which is $\pi/4$. Class AB sits just below. Real amplifiers reach the sixties.

\subsection{Why one Darlington is not the end of it}
\label{c8:sec:c2}
\secref{c8:sec:a5} reached 20 kilohm with two transistors. At the 3.5 A a 50 W amplifier needs, a power transistor's $h_{FE}$ is nearer 20:

\[
  I_{base} = \frac{3.54}{20 \times 20} = 8.8\ \text{mA}
\]

from a voltage-amplifier stage running at 1 to 10 mA. \textbf{That is the whole of its current, so the driver clips before the output does.} A \textbf{third} follower gives $h_{FE}^3$: the \textbf{triple emitter follower}, with a base current of 440 microamps.

The alternative is the \textbf{CFP-EF}, a complementary feedback pair driving an emitter follower. Local feedback round two devices makes the output track the input to within one $V_{BE}$ rather than two, and its thermal behaviour is set by the small driver rather than the hot power device, so it needs much less bias tracking (\secref{c8:sec:b4}). Its weakness is a tendency to oscillate. \textbf{Both exist for one reason:} current gain that does not depend on a hot power device's beta.

\subsection{The circuits that are only there to prevent failure}
\label{c8:sec:c3}
None of these affect the gain, and every one is in every commercial amplifier.

\begin{itemize}
  \item \textbf{Base stoppers,} 10 to 100 ohm in series with each output base. A power transistor and its own lead inductance oscillate at tens of megahertz, where nobody is looking. The resistor damps it.
  \item \textbf{The Zobel network,} typically 10 ohm and 100 nF from output to ground. A loudspeaker is inductive above a few kilohertz; the Zobel keeps something resistive there.
  \item \textbf{The output inductor,} a few microhenries wound over a resistor, isolating a capacitive cable whose phase lag would otherwise sit inside the feedback loop.
  \item \textbf{Snubbers} across the rectifier diodes. Reverse recovery steps current into the transformer's leakage inductance, which rings at radio frequency and reaches the output.
  \item \textbf{Over-current protection,} a transistor across each output base-emitter junction, turned on by the emitter resistor's voltage, stealing base drive above a set current. \textbf{So the emitter resistors of \secref{c8:sec:b3} do three jobs:} thermal stability, current sharing between paralleled devices, and current sensing.
  \item \textbf{A DC offset detector} and a relay in series with the loudspeaker. A shorted output device puts the full rail across 8 ohm, 128 W into a voice coil rated for a few.
\end{itemize}

\subsection{Heat}
\label{c8:sec:c4}
\textbf{Dissipation per device peaks at about 40 per cent of full output,} not at full output, where the current is high and the voltage across the device is still large. Here about 17 W: $V_{CC}^2/\pi^2R_L = 13$ W of class-B dissipation plus 4 W from the idle current.

\textbf{Thermal resistances add in series} like electrical ones: junction to case, case to heatsink through the washer, heatsink to air. 1 plus 0.5 plus 1 is 2.5 \textdegree{}C/W, so 17 W lifts the junction 42 \textdegree{}C: 82 \textdegree{}C at a 40 \textdegree{}C ambient, acceptable for a device rated to 150.

\textbf{Placement matters as much as size.} Every output device and the bias generator on the \emph{same} heatsink, close together, because \secref{c8:sec:b4} needs them at the same temperature. \textbf{Two devices on separate heatsinks will not share current, and the hotter one takes progressively more of it:} the same runaway one level up.

\subsection{Why output stages are usually bipolar}
\label{c8:sec:c5}
The opposite conclusion to the one \secref{c8:sec:b6} reaches for input stages:

\begin{booktable}{@{}llL@{}}
  \thead{} & \thead{Bipolar} & \thead{MOSFET} \\
  \midrule
  Output resistance & about ten times lower & higher, by $r_s/r_e$ \\
  Distortion & lower & higher without correction \\
  Thermal stability & needs tracking bias & inherently stable above a crossover point \\
  High-frequency stability & needs stoppers and Zobel & far less prone \\
  Drive & current, so beta matters & voltage, but gate charge matters \\
\end{booktable}

\textbf{The two devices swap places depending on which end of the amplifier they are at,} for one reason: a MOSFET's transconductance is ten times lower at the same current. At the input that is irrelevant and the infinite gate resistance decisive; at the output, the reverse. A MOSFET stage with an error amplifier can reach bipolar distortion figures, buying robustness rather than performance.

\subsection{What to take from this section}
\label{c8:sec:c6}
\begin{enumerate}
  \item \textbf{A power stage is mostly not the amplifier.} The gain path is a handful of transistors; the rest of the schematic is protection, compensation and thermal management.
  \item \textbf{The failures are thermal and high-frequency,} neither visible in the small-signal analysis this book teaches. A design correct in every equation of Chapter~\ref{c7:ch:smallsignal} and this chapter can still destroy itself.
  \item \textbf{The corrections in \secref{c8:sec:c1} are the kind that matter.} An rms mistaken for a peak is a rail 40 per cent too low, and nothing further down recovers it.
\end{enumerate}
